% TEMPLATE-MILC-v3.tex - plantilla maestra UNICA de las guias MILC v3.
%
% La usan las tres familias de guias, cada una con su driver:
%   · Grados 6-11  -> scripts/build-guias-g11.py
%   · Bebras       -> scripts/build-guias-bebras.py
%   · Semillero    -> scripts/build-guias-semillero.py
%
% Antes habia tres copias casi identicas (~400 lineas iguales, 6 distintas):
% cada mejora habia que aplicarla tres veces, y en la practica no se hacia
% (el motor de recursos vivia solo en dos de ellas). Lo que varia entre
% familias esta parametrizado en 6 placeholders que cada driver llena
% (se nombran aqui SIN su sintaxis real: el builder reemplaza por texto plano,
% tambien dentro de los comentarios, y un valor multilinea rompe el preambulo):
%
%   HEADER_IZQ        encabezado izquierdo de pagina
%   HEADER_DER        encabezado derecho de pagina
%   PORTADA_ETIQUETA  rotulo sobre el numero grande (GRADO/MOMENTO/MODULO)
%   PORTADA_SERIE     linea de serie de la portada
%   PORTADA_ID        identificador de la guia en la portada
%   PORTADA_CIRCULO   el circulito de grado (°), o vacio si no aplica
%
% El resto de claves las documenta content/guias/_SCHEMA.md. El script valida
% que no queden placeholders sin reemplazar antes de escribir el archivo final.

\documentclass[11pt,letterpaper]{article}
\usepackage[margin=1.75cm]{geometry}
\usepackage[table]{xcolor}
\usepackage{fontspec}
\usepackage[spanish]{babel}
\usepackage{tikz}
% Librerías TikZ para los diagramas de las guías (bloque `recursos:`):
%  · babel  → babel en español vuelve activos caracteres como «>», y TikZ los usa
%             en sus opciones (>=stealth, ->). Sin esta librería, cualquier
%             diagrama con flechas revienta con «\language@active@arg> has an
%             extra }». Debe ir ANTES de que se dibuje nada.
%  · positioning        → sintaxis «below left=2pt and 3pt».
%  · decorations.pathreplacing → llaves ({brace}) para señalar rangos.
\usetikzlibrary{babel,positioning,decorations.pathreplacing}
\usepackage{graphicx}
% Circuitos: el trabajo de electrónica se hace hoy con micro:bit y sus sensores
% integrados (MakeCode, sin cableado), así que no se necesitan esquemáticos.
% Si algún día entran montajes con protoboard, basta descomentar esta línea:
% los diagramas .tex/.tikz declarados en `recursos:` ya se incrustan solos.
% \usepackage{circuitikz}
\usepackage[most]{tcolorbox}
\usepackage{tabularx}
\usepackage{array}
\usepackage{fancyhdr}
\usepackage{titlesec}
\usepackage{ragged2e}
\usepackage{enumitem}
\usepackage{microtype}

% Helvetica Neue no trae la flecha U+2192, asi que cada «→» escrito literalmente
% en el YAML se compilaba a NADA, en silencio (462 flechas en 61 guias: rutas del
% tipo «paleta → tipografia → lockup» salian rotas). Mapeamos el caracter a su
% version matematica, que si tiene glifo. Asi el autor puede seguir escribiendo
% «→» comodamente en el YAML.
\usepackage{newunicodechar}
\newunicodechar{→}{\ensuremath{\rightarrow}}

\setmainfont{Helvetica Neue}
\setsansfont{Helvetica Neue}
\setlength{\parindent}{0pt}
\setlength{\parskip}{6pt}
\hyphenpenalty=200
\exhyphenpenalty=200
\pretolerance=400
\tolerance=2000
\setlength{\emergencystretch}{1em}
\renewcommand{\arraystretch}{1.25}
\setlist{leftmargin=5mm,itemsep=1mm,topsep=1mm}
\newcolumntype{Y}{>{\RaggedRight\arraybackslash}X}

\definecolor{milcMagenta}{HTML}{E6007E}
\definecolor{milcVino}{HTML}{5A0038}
\definecolor{milcTurquesa}{HTML}{0093A5}
\definecolor{milcVerde}{HTML}{58A923}
\definecolor{milcMostaza}{HTML}{E5B400}
\definecolor{milcOcre}{HTML}{B86600}
\definecolor{milcNegro}{HTML}{241816}
\definecolor{milcGris}{HTML}{F7F7F4}
\definecolor{milcLinea}{HTML}{E8E2DD}
\definecolor{milcGradePrimary}{HTML}{FF2D87}
\definecolor{milcGradeDark}{HTML}{9F1B56}
\definecolor{milcGradeSoft}{HTML}{FFE0EE}
\definecolor{milcGradeAccent}{HTML}{CFE7FF}
\definecolor{milcGradeText}{HTML}{FFFFFF}
\color{milcNegro}

\pagestyle{fancy}
\fancyhf{}
\lhead{\small Tecnología e Informática · Grado 10}
\rhead{\small Guía 7 · MILC v3}
\cfoot{\small\thepage}
\renewcommand{\headrulewidth}{0pt}

\newtcolorbox{titlebox}[1]{
  enhanced,colback=#1,colframe=#1,arc=7mm,boxrule=0pt,
  left=7mm,right=7mm,top=3mm,bottom=3mm,
  before skip=8pt,after skip=8pt,
  fontupper=\sffamily\bfseries\Large\color{white}
}

\newtcolorbox{softbox}[3]{
  enhanced,colback=#2,colframe=#1,borderline west={4pt}{0pt}{#1},
  arc=2mm,boxrule=.4pt,left=4mm,right=4mm,top=3mm,bottom=3mm,
  before skip=6pt,after skip=8pt,title={#3},coltitle=white,
  colbacktitle=#1,fonttitle=\sffamily\bfseries,fontupper=\RaggedRight,
  attach boxed title to top left={xshift=3mm,yshift=-2mm},
  boxed title style={arc=2mm, boxrule=0pt}
}

\newtcolorbox{drawbox}[2]{
  enhanced,colback=white,colframe=#1,arc=4mm,boxrule=1pt,
  left=5mm,right=5mm,top=4mm,bottom=4mm,before skip=8pt,after skip=8pt,
  title={#2},coltitle=white,colbacktitle=#1,fonttitle=\sffamily\bfseries,
  fontupper=\RaggedRight,
  attach boxed title to top left={xshift=4mm,yshift=-2mm},
  boxed title style={arc=3mm, boxrule=0pt}
}

\newtcolorbox{infoband}[1]{
  enhanced,colback=#1!8,colframe=#1!50,arc=2mm,boxrule=.5pt,
  left=4mm,right=4mm,top=2mm,bottom=2mm,before skip=4pt,after skip=4pt,
  fontupper=\small\RaggedRight
}

% Puente narrativo entre secciones (contrato editorial Conéctate).
% Da continuidad: "Acabas de X. Ahora vas a Y porque Z."
% Visualmente: banda izquierda en color del grado, texto en itálica pequeña.
\newtcolorbox{puentebox}{
  enhanced,colback=milcGradeSoft,colframe=milcGradeDark,
  borderline west={3pt}{0pt}{milcGradeDark},
  arc=0mm,boxrule=0pt,left=5mm,right=4mm,top=2.5mm,bottom=2.5mm,
  before skip=4pt,after skip=8pt,
  fontupper=\itshape\small\color{milcNegro}\RaggedRight
}
% La flecha va en modo matematico a proposito: Helvetica Neue NO tiene el glifo
% U+2192, asi que \textrightarrow se compilaba a NADA («Missing character: There
% is no → in font Helvetica Neue») y el puente salia sin flecha. En modo
% matematico la flecha viene de la fuente matematica y si se dibuja.
\newcommand{\puente}[1]{%
  \begin{puentebox}%
    \textcolor{milcGradeDark}{$\rightarrow$}\hspace{2mm}#1%
  \end{puentebox}%
}

\newcommand{\linead}[1]{\noindent\color{milcLinea}\rule{#1}{0.5pt}\color{milcNegro}}
\newcommand{\respuesta}[1]{\par\vspace{2mm}\foreach \n in {1,...,#1}{\noindent\color{milcLinea}\rule{\linewidth}{0.5pt}\par\vspace{5mm}}\color{milcNegro}}
\newcommand{\checkbox}{\textcolor{milcGradeDark}{\fbox{\phantom{x}}}\hspace{5pt}}

\newcommand{\phasecard}[4]{
\begin{tcolorbox}[enhanced,colback=#3,colframe=#2,arc=3mm,boxrule=.5pt,height=3.05cm,valign=center,halign=center,left=1.5mm,right=1.5mm,top=2mm,bottom=2mm]
{\sffamily\bfseries\fontsize{22}{24}\selectfont\textcolor{#2}{#1}}\par
{\sffamily\bfseries\small #4}
\end{tcolorbox}
}

% ── Recursos visuales de la guía ──────────────────────────────────────
% Declarados en el bloque `recursos:` del YAML y emitidos por el builder.
% \guiaFigura   → imagen rasterizada o PDF (foto, carta celeste, montaje).
% \guiaDiagrama → fuente TikZ/circuitikz (.tex) incrustada como vector:
%                 es la vía para circuitos y esquemáticos editables.
% #1 = ruta relativa al .tex · #2 = pie (puede ir vacío)
\newcommand{\guiaFigura}[2]{%
\begin{center}
\includegraphics[width=0.88\linewidth,height=0.38\textheight,keepaspectratio]{#1}\\[1.5mm]
{\footnotesize\itshape\textcolor{milcNegro!75}{#2}}
\end{center}
}

\newcommand{\guiaDiagrama}[2]{%
\begin{center}
\input{#1}\\[1.5mm]
{\footnotesize\itshape\textcolor{milcNegro!75}{#2}}
\end{center}
}

\begin{document}
\thispagestyle{empty}
\begin{tikzpicture}[remember picture,overlay]
  \fill[white] (current page.south west) rectangle (current page.north east);
  \path[fill=milcGradePrimary,rounded corners=16mm]
    ([xshift=.55cm,yshift=-.55cm]current page.north west) rectangle
    ([xshift=-.55cm,yshift=3.65cm]current page.south east);

  \node[anchor=north west,text=milcGradeText] at ([xshift=2.0cm,yshift=-1.85cm]current page.north west)
    {\sffamily\bfseries\fontsize{14}{16}\selectfont GRADO};
  \node[anchor=north west,text=milcGradeText,opacity=.82] at ([xshift=2.0cm,yshift=-2.75cm]current page.north west)
    {\sffamily\bfseries\fontsize{9}{11}\selectfont SERIE GUÍAS · TECNOLOGÍA E INFORMÁTICA};

  \begin{scope}[shift={([xshift=-3.05cm,yshift=-2.55cm]current page.north east)}]
    \fill[white,opacity=.80] (-.62,-.52) rectangle (.62,.52);
    \draw[milcGradeText,line width=1pt] (-.62,-.52) rectangle (.62,.52);
    \fill[milcGradeDark] (-.42,-.38) rectangle (-.22,.06);
    \fill[milcGradeAccent] (-.10,-.38) rectangle (.10,.24);
    \fill[milcGradePrimary!60!black] (.22,-.38) rectangle (.42,.38);
  \end{scope}

  \node[anchor=west,text=milcGradeText] at ([xshift=1.9cm,yshift=-12.85cm]current page.north west)
    {\sffamily\bfseries\fontsize{124}{124}\selectfont 10};
  % El circulito de grado (°) solo aplica a las guias de grado («11°»); en
  % Bebras y Semillero el numero grande es un momento/modulo, no un grado.
  \draw[milcGradeText,line width=1.7pt]
    ([xshift=4.52cm,yshift=-11.68cm]current page.north west) circle (.13cm);

  \node[anchor=west,text=milcGradeText,text width=8.7cm,align=left] at ([xshift=10.35cm,yshift=-11.6cm]current page.north west)
    {
     {\sffamily\bfseries\fontsize{19}{22}\selectfont Portada con IA ---\\herramientas gratuitas\\para generar imagen}\\[4mm]
     {\sffamily\bfseries\fontsize{23}{26}\selectfont Guía 7-10-TIC}\\[2mm]
     {\sffamily\fontsize{13}{16}\selectfont Período 1 · Escritura abierta con IA}\\[5mm]
     {\sffamily\bfseries\fontsize{11}{13}\selectfont Institución Educativa Sor María Juliana}\\[1mm]
     {\sffamily\fontsize{10}{12}\selectfont Autor: Dr. Álvaro Cárdenas Orozco}
    };

  \path[fill=milcGradeSoft,rounded corners=7mm]
    ([xshift=1.35cm,yshift=.85cm]current page.south west) rectangle
    ([xshift=-1.35cm,yshift=4.25cm]current page.south east);
  \draw[milcGradeDark,line width=.65pt,rounded corners=7mm]
    ([xshift=1.35cm,yshift=.85cm]current page.south west) rectangle
    ([xshift=-1.35cm,yshift=4.25cm]current page.south east);
  \node[anchor=center,text=milcNegro,text width=17.0cm,align=center] at ([yshift=2.55cm]current page.south)
    {\sffamily\fontsize{10}{12}\selectfont Metodología MILC · Apertura ancestral · Pensamiento computacional · Triángulo Dussel-Estoicismo-Floridi\\
    \textcolor{milcGradeDark}{\bfseries Producto:} Portada definitiva de tu libro creada con al menos 1 herramienta gratuita de IA generativa de imagen (Bing Image Creator, Adobe Express, Leonardo.ai, Microsoft Designer, ideogram.ai) + 3 versiones probadas con prompts distintos + decisión final justificada por escrito + composición final agregando título del libro y nombre del autor sobre la imagen + declaración del modelo usado y respeto a los términos de uso de la herramienta};
\end{tikzpicture}
\mbox{}
\newpage

\section*{Apertura: Portada con IA generativa de imagen --- herramientas gratuitas}

\begin{softbox}{milcGradeDark}{milcGradeSoft}{Saber ancestral}
En la carrera 5 del centro de Cartago, en los talleres de pintura de carteles que sostuvieron la comunicación visual del pueblo antes de las imprentas digitales, había un personaje que cualquier comerciante del barrio recuerda: \textbf{el pintor de carteles}. Cuando alguien quería anunciar tienda, velorio, fiesta patronal o publicación nueva, iba al pintor con la idea. El pintor nunca tomaba el pincel y empezaba. Tenía un ritual previo silencioso: \textbf{bocetaba 3 versiones a lápiz}. Primera versión: idea visual obvia. Segunda: variación con cambio de composición. Tercera: algo distinto, atrevido, inesperado. Ponía los 3 bocetos sobre la mesa, llamaba al cliente, pedía que escogiera. El cliente miraba, comentaba, a veces decía \emph{``mezcla la 1 con la 3''}. Solo después de la elección, el pintor tomaba el pincel. La sabiduría era simple: \textbf{una sola idea visual rara vez es la mejor; 3 versiones permiten elegir con criterio}. La portada con IA generativa es ese boceto multiplicado por 100. Pero el oficio del editor (elegir, decidir, firmar) sigue siendo humano.
\end{softbox}

\begin{softbox}{milcTurquesa}{milcGris}{Saber contemporáneo}
Las \textbf{herramientas gratuitas de IA generativa de imagen} en 2026 son varias, cada una con fortalezas distintas: (1) \textbf{Bing Image Creator} (\texttt{bing.com/images/create}): basado en DALL-E 3, gratis, requiere cuenta Microsoft. Bueno para estilos variados. (2) \textbf{Adobe Express} (\texttt{adobe.com/express}): incluye \emph{text to image} en plan gratuito limitado. (3) \textbf{Leonardo.ai}: 150 imágenes gratis al día con cuenta. Mucho control sobre estilo. (4) \textbf{Microsoft Designer}: integra IA con composición de portadas, gratis con cuenta. (5) \textbf{Ideogram.ai}: especializado en imágenes con texto integrado, gratis con límite diario. La regla profesional: \textbf{cada herramienta tiene su estética; conviene probar 2-3 antes de decidir}. La estructura del \textbf{prompt para imagen} es distinta del prompt para texto: importan más los descriptores visuales (estilo artístico, paleta de color, composición, atmósfera) que las instrucciones narrativas. Términos clave: \emph{realistic}, \emph{watercolor}, \emph{cinematic}, \emph{minimalist}, \emph{vibrant}, \emph{moody}. Cada herramienta tiene términos de uso propios que conviene leer antes de comercializar el resultado.
\end{softbox}

\begin{softbox}{milcMostaza}{milcGris}{Pregunta-puente}
\textit{¿Qué sabía el pintor de carteles al bocetar 3 versiones antes de pintar la definitiva, que el usuario novato de DALL-E olvida cuando acepta la primera imagen que sale? ¿Y por qué el prompt para imagen requiere descriptores visuales en lugar de instrucciones narrativas?}
\end{softbox}

\begin{softbox}{milcVerde}{milcGris}{Saber y saber hacer}
Al terminar podrás: (1) \textbf{identificar} las 3-4 herramientas gratuitas de IA generativa de imagen más útiles y sus diferencias de estilo; (2) \textbf{aplicar} los descriptores visuales correctos para construir prompts de imagen efectivos; (3) \textbf{crear} 3 versiones de portada con prompts distintos y elegir la definitiva con criterio editorial; (4) \textbf{evaluar} la portada final agregando título y autor en composición final.
\end{softbox}



\small
\begin{tabularx}{\textwidth}{>{\bfseries}p{3.0cm}Y}
\rowcolor{milcGradePrimary!12}Autor & Dr. Álvaro Cárdenas Orozco\\
DBA & Producción editorial mediada por IA con criterio ético y técnico (MEN, T\&I)\\
Referentes & Ley 23 de 1982 · Floridi (ética de la IA generativa) · Dussel (autoría con dignidad)\\
Producto final & Portada definitiva de tu libro creada con al menos 1 herramienta gratuita de IA generativa de imagen (Bing Image Creator, Adobe Express, Leonardo.ai, Microsoft Designer, ideogram.ai) + 3 versiones probadas con prompts distintos + decisión final justificada por escrito + composición final agregando título del libro y nombre del autor sobre la imagen + declaración del modelo usado y respeto a los términos de uso de la herramienta\\
\end{tabularx}
\normalsize

\puente{Antes de empezar, mira el viaje completo. Tienes texto en V3, estructura cerrada, derechos definidos. Hoy le pones la \textbf{primera imagen pública} a tu libro: la portada que verá quien lo encuentre. Las próximas 2 sesiones (diagramación, producción) integrarán la portada con el contenido.}

\begin{titlebox}{milcGradeDark}
Ruta MILC: así vamos a aprender
\end{titlebox}
\begin{tabularx}{\textwidth}{>{\centering\arraybackslash}X>{\centering\arraybackslash}X>{\centering\arraybackslash}X>{\centering\arraybackslash}X}
\phasecard{1}{milcVerde}{milcGris}{Escucha\\[-1mm]\small Observo, escucho y pregunto.} &
\phasecard{2}{milcTurquesa}{milcGris}{Sistematización\\[-1mm]\small Pensamiento computacional.} &
\phasecard{3}{milcMagenta}{milcGris}{Praxis\\[-1mm]\small Construyo y aplico.} &
\phasecard{4}{milcVino}{milcGris}{Evaluación\\[-1mm]\small Reflexión liberadora.}
\end{tabularx}


\puente{Antes de teorizar, vas a hacer una \emph{auditoría rápida} de portadas: piensa en 3 portadas de libros que recuerdas (de cualquier libro que conozcas) y anota qué imagen tienen, qué colores, qué tipografía, qué te llamó la atención cuando las viste por primera vez.}

\begin{titlebox}{milcVerde}
Fase 1 · Escucha
\end{titlebox}

\begin{softbox}{milcVerde}{milcGris}{Escena para escuchar}
\textbf{Actividad 1 · IDENTIFICA --- 3 portadas referencia} (15 min · individual). Piensa en 3 portadas de libros que recuerdas vívidamente (puede ser un libro escolar, una novela, un manga, un libro de YouTuber). Para cada una, anota: (1) ¿qué imagen tiene (persona, objeto, abstracta, paisaje)? (2) ¿qué paleta de colores domina? (3) ¿cómo aparece el título (grande arriba, integrado a la imagen, abajo)? (4) ¿qué te hizo recordarla (impacto, simplicidad, color, misterio)? Esa observación intuitiva te enseña qué hace funcionar una portada.
\end{softbox}

\checkbox Elegí 3 portadas que recuerdo.\\
\checkbox Respondí 4 preguntas sobre cada una.\\
\checkbox Identifiqué qué las hace memorables.

\begin{infoband}{milcVerde}


\textbf{Lo que va al cuaderno (Actividad 1 · IDENTIFICA).} Título: «Actividad 1 --- 3 portadas referencia». \textbf{Formato:} ficha por cada portada con imagen, color, tipografía, impacto. \textbf{Sabes que terminaste cuando} ves qué elementos comunes hacen funcionar a una buena portada.
\end{infoband}

\puente{Ya tienes 3 portadas referencia. Ahora vas a entender la \textbf{anatomía del prompt para imagen}: descriptores visuales, estilos, atmósfera, los 5 errores típicos (prompt narrativo en lugar de visual, esperar texto perfecto en imagen, ignorar términos de uso, aceptar la primera imagen, sin composición final con título).}

\begin{titlebox}{milcTurquesa}
Fase 2 · Sistematización con pensamiento computacional
\end{titlebox}

\begin{softbox}{milcTurquesa}{milcGris}{Cuatro pilares aplicados al tema}
Tus 3 referencias te dan vocabulario visual. Ahora necesitas la \textbf{anatomía profesional} del prompt para imagen y las herramientas disponibles.
\end{softbox}

\small
\begin{tabularx}{\textwidth}{>{\bfseries\RaggedRight\arraybackslash}p{3.6cm}Y}
\rowcolor{milcTurquesa!90}\color{white}Pilar & \color{white}\bfseries Aplicación al tema\\
1. Descomposición & Los \textbf{descriptores visuales} más útiles para el prompt de imagen se descomponen así: (1) \textbf{Sujeto principal}: qué se ve (\emph{a teenage girl with headphones}, \emph{a tropical bird in a forest}, \emph{an open book on a wooden desk}). (2) \textbf{Estilo artístico}: \emph{photorealistic}, \emph{watercolor}, \emph{digital illustration}, \emph{minimalist vector}, \emph{oil painting}, \emph{anime style}, \emph{pencil sketch}. (3) \textbf{Paleta y atmósfera}: \emph{vibrant warm colors}, \emph{moody dark tones}, \emph{pastel palette}, \emph{high contrast black and white}, \emph{golden hour lighting}. (4) \textbf{Composición}: \emph{centered}, \emph{rule of thirds}, \emph{close-up}, \emph{wide shot}, \emph{symmetrical}. (5) \textbf{Detalles contextuales}: ambiente, objetos secundarios, época, lugar.\\
2. Reconocimiento de patrones & Las \textbf{4 herramientas gratuitas} principales en 2026: (a) \textbf{Bing Image Creator}: rápido, basado en DALL-E 3, generoso con créditos gratuitos. Excelente como primera opción. (b) \textbf{Leonardo.ai}: 150 imágenes/día gratis, mucho control sobre modelos y estilos. (c) \textbf{Microsoft Designer}: bueno para combinar imagen IA + composición final con texto y plantillas. (d) \textbf{Ideogram.ai}: el mejor para integrar texto legible dentro de la imagen (otros modelos fallan con texto).\\
3. Abstracción & Lo esencial cabe en una pregunta: \emph{¿esta imagen captura la esencia del libro?}. Si tu libro habla de adolescentes y ansiedad, una imagen genérica de \emph{``persona sonriendo''} no captura. Si tu libro es manual de emprendimiento, una imagen \emph{``corporativa estándar''} tampoco. La phronesis del editor consiste en \textbf{rechazar imágenes que no representan el contenido, aunque sean técnicamente buenas}.\\
4. Algoritmo conceptual & Algoritmo: (1) revisar tus 3 portadas referencia. (2) escribir 3 prompts distintos (3 enfoques visuales para tu libro). (3) probar cada prompt en al menos 1 herramienta gratuita. (4) generar al menos 3-4 imágenes por prompt. (5) elegir las 3 mejores en total. (6) decidir la definitiva con criterio. (7) agregar título y autor en composición final usando Canva o Microsoft Designer.\\
\end{tabularx}
\normalsize

\begin{drawbox}{milcTurquesa}{Anatomía del prompt para imagen}
\textbf{Sujeto:} qué se ve, en 3-5 palabras. \\[1mm] \textbf{Estilo:} artístico (foto, acuarela, digital, etc.). \\[1mm] \textbf{Paleta:} colores y atmósfera. \\[1mm] \textbf{Composición:} encuadre y enfoque. \\[1mm] \textbf{Detalles:} ambiente y contexto.
\end{drawbox}

\begin{softbox}{milcMostaza}{milcGris}{Errores comunes y su corrección}
(1) \textbf{Prompt narrativo en vez de visual}: \emph{``una niña que está triste por algo''} en lugar de \emph{``close-up of a teenage girl, dim lighting, watercolor style, melancholic mood''}. (2) \textbf{Esperar texto perfecto en imagen}: la mayoría de modelos fallan con texto integrado (excepto Ideogram). (3) \textbf{Ignorar términos de uso}: cada herramienta tiene reglas distintas sobre uso comercial. (4) \textbf{Aceptar la primera imagen}: similar al primer borrador, casi nunca es la mejor. (5) \textbf{Sin composición final con título}: la imagen IA es solo la base; el título y autor se agregan después con Canva o similar.
\end{softbox}

\begin{infoband}{milcTurquesa}


\textbf{Lo que va al cuaderno (Actividad 2 · APLICA).} Título: «Actividad 2 --- Anatomía prompt imagen». \textbf{Formato:} (a) los 5 descriptores con ejemplo; (b) 4 herramientas con su uso; (c) los 5 errores con marca. \textbf{Sabes que terminaste cuando} podrías escribir un prompt visual sin abrir esta guía.
\end{infoband}

\puente{Tienes el método. Toca aplicarlo: escribes 3 prompts visuales distintos, los pruebas en 1-2 herramientas, eliges la versión definitiva, agregas título y autor en composición final.}

\begin{titlebox}{milcMagenta}
Fase 3 · Praxis con pensamiento computacional
\end{titlebox}

\begin{softbox}{milcMagenta}{milcGris}{Construyendo el producto con los cuatro pilares}
\textbf{Actividad 3 · CREA + EVALÚA --- 3 versiones + portada definitiva + composición final} (40 min · individual con dispositivo). Hora de aplicar. Escribes 3 prompts visuales distintos, generas imágenes, eliges, compones con título y autor.
\end{softbox}

\small
\begin{tabularx}{\textwidth}{>{\bfseries\RaggedRight\arraybackslash}p{3.6cm}Y}
\rowcolor{milcMagenta!90}\color{white}Pilar & \color{white}\bfseries Aplicación al producto\\
Descomposición & Tu portada se construye en 6 pasos: (1) escribir 3 prompts visuales distintos para tu libro (3 enfoques: realista, ilustrado, conceptual). (2) probar cada uno en al menos 1 herramienta (Bing Image Creator es buen primer paso). (3) generar 3-4 imágenes por prompt. (4) elegir las 3 mejores en total. (5) decidir la versión definitiva con criterio editorial (¿cuál representa mejor el libro?). (6) llevar la imagen elegida a Canva o Microsoft Designer y agregar título grande, autor más pequeño, en composición final.\\
Patrones & Sigue el patrón \textbf{``imagen IA + composición humana''}: la imagen pura de la IA rara vez es portada final. Necesita agregarse título legible, nombre del autor, posiblemente subtítulo. Esa composición es trabajo humano del editor con herramientas gratuitas (Canva, Designer, Express). La portada final es la unión de imagen generada con composición curada.\\
Abstracción & La decisión justificada por escrito tiene 3 puntos: (a) cuál elegí, (b) por qué representa mejor el libro, (c) qué descarté de las otras dos y por qué. 1 párrafo total. Esta justificación demuestra criterio editorial, no aceptación pasiva de la IA.\\
Algoritmo de armado & Algoritmo: (1) escribir 3 prompts. (2) generar imágenes. (3) elegir 3 mejores. (4) decidir definitiva. (5) componer con título y autor. (6) justificar decisión. (7) declarar modelo usado y respetar términos de uso de la herramienta.\\
\end{tabularx}
\normalsize

\begin{drawbox}{milcMagenta}{Checklist antes de entregar}
\checkbox 3 prompts visuales distintos escritos. \\[1mm] \checkbox Probados en al menos 1 herramienta gratuita. \\[1mm] \checkbox 3 versiones de portada generadas. \\[1mm] \checkbox Decisión final justificada por escrito. \\[1mm] \checkbox Composición final con título y autor agregados. \\[1mm] \checkbox Modelo de IA declarado por escrito. \\[1mm] \checkbox Términos de uso de la herramienta verificados.
\end{drawbox}

\begin{softbox}{milcMagenta}{milcGris}{Plantilla de trabajo}
\textbf{Prompt 1.} \textit{Enfoque realista o fotográfico.} \\[1mm] \textbf{Prompt 2.} \textit{Enfoque ilustrado o artístico.} \\[1mm] \textbf{Prompt 3.} \textit{Enfoque conceptual o abstracto.} \\[1mm] \textbf{Decisión.} \textit{Cuál elegí + por qué + qué descarté.} \\[1mm] \textbf{Composición.} \textit{Imagen + título + autor en herramienta gratuita.} \\[1mm] \textbf{Declaración.} \textit{Modelo IA usado + términos de uso.}
\end{softbox}

\begin{infoband}{milcMagenta}


\textbf{Lo que va al cuaderno (Actividad 3 · CREA + EVALÚA).} Título: «Actividad 3 --- Mi portada». \textbf{Formato:} (a) los 3 prompts; (b) referencias a las 3 versiones generadas; (c) portada definitiva con composición; (d) decisión justificada; (e) declaración del modelo. \textbf{Sabes que terminaste cuando} tienes portada lista para diagramación.
\end{infoband}

\puente{Lo que sigue resume \emph{qué} entregas hoy: 3 versiones + decisión + composición final + declaración del modelo. El criterio principal: que cualquiera viendo la portada entienda qué libro es y se sienta llamado a abrirlo.}

\begin{titlebox}{milcMostaza}
Producto de la sesión
\end{titlebox}
\textbf{3 versiones + portada definitiva + composición final + declaración del modelo}.
\begin{softbox}{milcMostaza}{milcGris}{Criterios de calidad}
(1) 3 prompts visuales distintos. (2) Probados en herramientas gratuitas. (3) 3 versiones de portada generadas. (4) Decisión justificada con criterio editorial. (5) Composición final con título y autor. (6) Declaración del modelo usado.
\end{softbox}

\puente{Antes de cerrar, mira la portada desde las cinco dimensiones humanas. Elegir entre 3 versiones es disciplina del oficio; aceptar la primera es renunciar al criterio editorial.}

\begin{titlebox}{milcVino}
Fase 4 · Evaluación liberadora — cinco dimensiones
\end{titlebox}

\begin{softbox}{milcVino}{milcGris}{Sentido de la evaluación}
Evaluar no es solo revisar una nota. Es reconocer qué aprendí, qué cambió en mí, qué emociones manejé, cómo me relaciono con la ciudad y la comunidad, y qué le dejaré a quien viene detrás.
\end{softbox}

\textbf{Mi reflexión liberadora en cinco dimensiones.} Completa cada frase iniciada:

\small
\begin{tabularx}{\textwidth}{>{\bfseries\RaggedRight\arraybackslash}p{3.4cm}Y>{\RaggedRight\arraybackslash}p{4.6cm}}
\rowcolor{milcVino}\color{white}Dimensión & \color{white}\bfseries Mi respuesta (frame starter) & \color{white}\bfseries Algo que descubrí\\
1. Personal & Lo que más cambió en mí fue \linead{6.5cm} & \linead{4.2cm}\\
2. Emocional & La emoción más fuerte fue \linead{2.5cm} y la regulé \linead{3cm} & \linead{4.2cm}\\
3. Ciudadana & En democracia esto importa porque \linead{6cm} & \linead{4.2cm}\\
4. Local (Cartago) & En mi barrio sería útil para \linead{6.5cm} & \linead{4.2cm}\\
5. Intergeneracional & A mi abuela/abuelo le diría \linead{6.5cm} & \linead{4.2cm}\\
\end{tabularx}
\normalsize

\begin{infoband}{milcVino}
\textbf{Para la próxima sustentación.} Vuelve a esta tabla en seis meses. Verás que las respuestas cambian — no porque lo aprendido cambie, sino porque tú cambias. La evaluación liberadora es una conversación contigo a través del tiempo.
\end{infoband}


\puente{Y para terminar, tres voces sobre la imagen del libro. Dussel sobre las portadas que dignifican al lector o lo manipulan, Marco Aurelio sobre la simplicidad como virtud visual, Floridi sobre la transparencia del autor de imagen IA.}

\begin{titlebox}{milcGradeDark}
Triángulo de pensamiento · cierre filosófico
\end{titlebox}

\begin{softbox}{milcVino}{milcGris}{Dussel · lente del nosotros}
\textit{``Una portada que representa al libro con dignidad respeta al lector; una que engaña visualmente lo manipula.''}

La portada es el primer contacto del lector con el libro. Si engaña sobre el contenido (imagen comercial que no representa el tema real), traiciona al lector. La filosofía de la liberación pide portadas que sean \textbf{anticipo honesto del contenido}, no carnada para venta. La phronesis del editor consiste en \textbf{elegir la portada que el lector aprobará después de leer, no la que más vende}.

\textbf{¿Mi portada anticipa con honestidad el contenido, o promete algo que el libro no cumple?}
\end{softbox}
\linead{\linewidth}\\[1mm]
\linead{\linewidth}

\begin{softbox}{milcMostaza}{milcGris}{Estoicismo · Marco Aurelio · cuidado interior}
\textit{``Lo simple comunica; lo recargado distrae. La virtud visual está en lo que se quita.''}

Es tentador llenar la portada con elementos para que se vea \emph{``profesional''}. La sabiduría estoica recomienda lo opuesto: \emph{quita hasta que solo quede lo esencial}. Las portadas más memorables del mundo tienen pocos elementos pero precisos. La phronesis visual del editor consiste en \textbf{resistir la tentación del adorno innecesario}.

\textbf{¿Mi portada tiene solo lo esencial, o agregué cosas que se pueden quitar sin perder mensaje?}
\end{softbox}
\linead{\linewidth}\\[1mm]
\linead{\linewidth}

\begin{softbox}{milcTurquesa}{milcGris}{Floridi · lente de la infoesfera}
\textit{``Declarar el uso de IA en la portada es transparencia editorial moderna; ocultarlo es fraude visual.''}

En la era de IA generativa, declarar que la portada fue creada con asistencia de IA es práctica ética emergente. No te resta autoría como editor: te suma transparencia profesional. Tu declaración del modelo usado es ya acto ético que muchos autores aún no hacen. Te entrena para una práctica que será estándar en pocos años.

\textbf{¿Mi declaración del modelo IA es transparente, o intenté ocultarlo para parecer más \emph{``auténtico''}?}
\end{softbox}
\linead{\linewidth}\\[1mm]
\linead{\linewidth}

\begin{titlebox}{milcVino}
Compromiso final
\end{titlebox}

\small
\begin{tabularx}{\textwidth}{>{\RaggedRight\arraybackslash}p{3.0cm}YYY}
\rowcolor{milcVino}\color{white}\bfseries Criterio & \color{white}\bfseries Lo logré & \color{white}\bfseries En proceso & \color{white}\bfseries Necesito apoyo\\
Comprensión del tema & Explico ideas centrales con mis palabras. & Reconozco ideas pero falta precisión. & Necesito repasar conceptos base.\\
Pensamiento computacional & Apliqué los 4 pilares al diseñar mi producto. & Apliqué algunos pilares. & Necesito releer la fase 2.\\
Aplicación y producto & Mi producto cumple los criterios y se puede revisar. & Producto funcional con ajustes pendientes. & Aún en construcción.\\
Reflexión 5 dimensiones & Respondí cada dimensión con honestidad. & Respondí algunas con detalle. & Mis respuestas son muy generales.\\
Triángulo de pensamiento & Conecté las 3 voces con mi trabajo real. & Conecté al menos una. & No relaciono las voces con mi proceso.\\
\end{tabularx}
\normalsize

\textbf{Hoy entendí que:}\\[1mm]
\linead{\linewidth}\\[1mm]
\linead{\linewidth}

\textbf{Mi compromiso para la próxima sesión es:}\\[1mm]
\linead{\linewidth}\\[1mm]
\linead{\linewidth}

\begin{softbox}{milcMostaza}{milcGris}{Cita de despedida · Marco Aurelio}
\textit{``El obstáculo en el camino se vuelve el camino. Lo que estorba al camino se vuelve camino.''}\\[1mm]
Cada error de hoy, bien observado, se convierte en pieza del camino que pisarás mañana.
\end{softbox}

\small
\begin{tabularx}{\textwidth}{>{\bfseries}p{3.6cm}Y}
\rowcolor{milcGradeDark!12}Nombre completo: & \linead{10cm}\\
Fecha: & \linead{4cm} \quad Curso/grupo: \linead{4cm}\\
Mi firma: & \linead{9cm}\\
\end{tabularx}
\normalsize

\begin{infoband}{milcGradeDark}
\textbf{Para la próxima sesión.} Llega con tu portada definitiva, las 3 versiones probadas y la declaración del modelo. En la sesión 8 vas a aprender \textbf{diagramación con herramientas gratuitas}: Canva, Google Docs, LibreOffice. La portada de hoy es el rostro; la diagramación de mañana es el cuerpo del libro.
\end{infoband}

\end{document}
